%% ODER: format ==         = "\mathrel{==}"
%% ODER: format /=         = "\neq "
%
%
\makeatletter
\@ifundefined{lhs2tex.lhs2tex.sty.read}%
  {\@namedef{lhs2tex.lhs2tex.sty.read}{}%
   \newcommand\SkipToFmtEnd{}%
   \newcommand\EndFmtInput{}%
   \long\def\SkipToFmtEnd#1\EndFmtInput{}%
  }\SkipToFmtEnd

\newcommand\ReadOnlyOnce[1]{\@ifundefined{#1}{\@namedef{#1}{}}\SkipToFmtEnd}
\usepackage{amstext}
\usepackage{amssymb}
\usepackage{stmaryrd}
\DeclareFontFamily{OT1}{cmtex}{}
\DeclareFontShape{OT1}{cmtex}{m}{n}
  {<5><6><7><8>cmtex8
   <9>cmtex9
   <10><10.95><12><14.4><17.28><20.74><24.88>cmtex10}{}
\DeclareFontShape{OT1}{cmtex}{m}{it}
  {<-> ssub * cmtt/m/it}{}
\newcommand{\texfamily}{\fontfamily{cmtex}\selectfont}
\DeclareFontShape{OT1}{cmtt}{bx}{n}
  {<5><6><7><8>cmtt8
   <9>cmbtt9
   <10><10.95><12><14.4><17.28><20.74><24.88>cmbtt10}{}
\DeclareFontShape{OT1}{cmtex}{bx}{n}
  {<-> ssub * cmtt/bx/n}{}
\newcommand{\tex}[1]{\text{\texfamily#1}}	% NEU

\newcommand{\Sp}{\hskip.33334em\relax}


\newcommand{\Conid}[1]{\mathit{#1}}
\newcommand{\Varid}[1]{\mathit{#1}}
\newcommand{\anonymous}{\kern0.06em \vbox{\hrule\@width.5em}}
\newcommand{\plus}{\mathbin{+\!\!\!+}}
\newcommand{\bind}{\mathbin{>\!\!\!>\mkern-6.7mu=}}
\newcommand{\rbind}{\mathbin{=\mkern-6.7mu<\!\!\!<}}% suggested by Neil Mitchell
\newcommand{\sequ}{\mathbin{>\!\!\!>}}
\renewcommand{\leq}{\leqslant}
\renewcommand{\geq}{\geqslant}
\usepackage{polytable}

%mathindent has to be defined
\@ifundefined{mathindent}%
  {\newdimen\mathindent\mathindent\leftmargini}%
  {}%

\def\resethooks{%
  \global\let\SaveRestoreHook\empty
  \global\let\ColumnHook\empty}
\newcommand*{\savecolumns}[1][default]%
  {\g@addto@macro\SaveRestoreHook{\savecolumns[#1]}}
\newcommand*{\restorecolumns}[1][default]%
  {\g@addto@macro\SaveRestoreHook{\restorecolumns[#1]}}
\newcommand*{\aligncolumn}[2]%
  {\g@addto@macro\ColumnHook{\column{#1}{#2}}}

\resethooks

\newcommand{\onelinecommentchars}{\quad-{}- }
\newcommand{\commentbeginchars}{\enskip\{-}
\newcommand{\commentendchars}{-\}\enskip}

\newcommand{\visiblecomments}{%
  \let\onelinecomment=\onelinecommentchars
  \let\commentbegin=\commentbeginchars
  \let\commentend=\commentendchars}

\newcommand{\invisiblecomments}{%
  \let\onelinecomment=\empty
  \let\commentbegin=\empty
  \let\commentend=\empty}

\visiblecomments

\newlength{\blanklineskip}
\setlength{\blanklineskip}{0.66084ex}

\newcommand{\hsindent}[1]{\quad}% default is fixed indentation
\let\hspre\empty
\let\hspost\empty
\newcommand{\NB}{\textbf{NB}}
\newcommand{\Todo}[1]{$\langle$\textbf{To do:}~#1$\rangle$}

\EndFmtInput
\makeatother
%
%
%
%
%
%
% This package provides two environments suitable to take the place
% of hscode, called "plainhscode" and "arrayhscode". 
%
% The plain environment surrounds each code block by vertical space,
% and it uses \abovedisplayskip and \belowdisplayskip to get spacing
% similar to formulas. Note that if these dimensions are changed,
% the spacing around displayed math formulas changes as well.
% All code is indented using \leftskip.
%
% Changed 19.08.2004 to reflect changes in colorcode. Should work with
% CodeGroup.sty.
%
\ReadOnlyOnce{polycode.fmt}%
\makeatletter

\newcommand{\hsnewpar}[1]%
  {{\parskip=0pt\parindent=0pt\par\vskip #1\noindent}}

% can be used, for instance, to redefine the code size, by setting the
% command to \small or something alike
\newcommand{\hscodestyle}{}

% The command \sethscode can be used to switch the code formatting
% behaviour by mapping the hscode environment in the subst directive
% to a new LaTeX environment.

\newcommand{\sethscode}[1]%
  {\expandafter\let\expandafter\hscode\csname #1\endcsname
   \expandafter\let\expandafter\endhscode\csname end#1\endcsname}

% "compatibility" mode restores the non-polycode.fmt layout.

\newenvironment{compathscode}%
  {\par\noindent
   \advance\leftskip\mathindent
   \hscodestyle
   \let\\=\@normalcr
   \let\hspre\(\let\hspost\)%
   \pboxed}%
  {\endpboxed\)%
   \par\noindent
   \ignorespacesafterend}

\newcommand{\compaths}{\sethscode{compathscode}}

% "plain" mode is the proposed default.
% It should now work with \centering.
% This required some changes. The old version
% is still available for reference as oldplainhscode.

\newenvironment{plainhscode}%
  {\hsnewpar\abovedisplayskip
   \advance\leftskip\mathindent
   \hscodestyle
   \let\hspre\(\let\hspost\)%
   \pboxed}%
  {\endpboxed%
   \hsnewpar\belowdisplayskip
   \ignorespacesafterend}

\newenvironment{oldplainhscode}%
  {\hsnewpar\abovedisplayskip
   \advance\leftskip\mathindent
   \hscodestyle
   \let\\=\@normalcr
   \(\pboxed}%
  {\endpboxed\)%
   \hsnewpar\belowdisplayskip
   \ignorespacesafterend}

% Here, we make plainhscode the default environment.

\newcommand{\plainhs}{\sethscode{plainhscode}}
\newcommand{\oldplainhs}{\sethscode{oldplainhscode}}
\plainhs

% The arrayhscode is like plain, but makes use of polytable's
% parray environment which disallows page breaks in code blocks.

\newenvironment{arrayhscode}%
  {\hsnewpar\abovedisplayskip
   \advance\leftskip\mathindent
   \hscodestyle
   \let\\=\@normalcr
   \(\parray}%
  {\endparray\)%
   \hsnewpar\belowdisplayskip
   \ignorespacesafterend}

\newcommand{\arrayhs}{\sethscode{arrayhscode}}

% The mathhscode environment also makes use of polytable's parray 
% environment. It is supposed to be used only inside math mode 
% (I used it to typeset the type rules in my thesis).

\newenvironment{mathhscode}%
  {\parray}{\endparray}

\newcommand{\mathhs}{\sethscode{mathhscode}}

% texths is similar to mathhs, but works in text mode.

\newenvironment{texthscode}%
  {\(\parray}{\endparray\)}

\newcommand{\texths}{\sethscode{texthscode}}

% The framed environment places code in a framed box.

\def\codeframewidth{\arrayrulewidth}
\RequirePackage{calc}

\newenvironment{framedhscode}%
  {\parskip=\abovedisplayskip\par\noindent
   \hscodestyle
   \arrayrulewidth=\codeframewidth
   \tabular{@{}|p{\linewidth-2\arraycolsep-2\arrayrulewidth-2pt}|@{}}%
   \hline\framedhslinecorrect\\{-1.5ex}%
   \let\endoflinesave=\\
   \let\\=\@normalcr
   \(\pboxed}%
  {\endpboxed\)%
   \framedhslinecorrect\endoflinesave{.5ex}\hline
   \endtabular
   \parskip=\belowdisplayskip\par\noindent
   \ignorespacesafterend}

\newcommand{\framedhslinecorrect}[2]%
  {#1[#2]}

\newcommand{\framedhs}{\sethscode{framedhscode}}

% The inlinehscode environment is an experimental environment
% that can be used to typeset displayed code inline.

\newenvironment{inlinehscode}%
  {\(\def\column##1##2{}%
   \let\>\undefined\let\<\undefined\let\\\undefined
   \newcommand\>[1][]{}\newcommand\<[1][]{}\newcommand\\[1][]{}%
   \def\fromto##1##2##3{##3}%
   \def\nextline{}}{\) }%

\newcommand{\inlinehs}{\sethscode{inlinehscode}}

% The joincode environment is a separate environment that
% can be used to surround and thereby connect multiple code
% blocks.

\newenvironment{joincode}%
  {\let\orighscode=\hscode
   \let\origendhscode=\endhscode
   \def\endhscode{\def\hscode{\endgroup\def\@currenvir{hscode}\\}\begingroup}
   %\let\SaveRestoreHook=\empty
   %\let\ColumnHook=\empty
   %\let\resethooks=\empty
   \orighscode\def\hscode{\endgroup\def\@currenvir{hscode}}}%
  {\origendhscode
   \global\let\hscode=\orighscode
   \global\let\endhscode=\origendhscode}%

\makeatother
\EndFmtInput
%

% format { = "\lbag"
% format } = "\rbag"


% TODO: fix spacing for \lbag and \rbag





















































%formap phiX = "\Phi"














\section{A quick tour of BiGUL}

\ignore{


\begin{hscode}\SaveRestoreHook
\column{B}{@{}>{\hspre}l<{\hspost}@{}}%
\column{E}{@{}>{\hspre}l<{\hspost}@{}}%
\>[B]{}\mbox{\enskip\{-\# LANGUAGE FlexibleContexts, TemplateHaskell, TypeFamilies  \#-\}\enskip}{}\<[E]%
\\[\blanklineskip]%
\>[B]{}\mathbf{module}\;\Conid{PBasic}\;\mathbf{where}{}\<[E]%
\\
\>[B]{}\mathbf{import}\;\Conid{\Conid{Generics}.BiGUL}{}\<[E]%
\\
\>[B]{}\mathbf{import}\;\Conid{\Conid{Generics}.\Conid{BiGUL}.Interpreter}{}\<[E]%
\\
\>[B]{}\mathbf{import}\;\Conid{\Conid{Generics}.\Conid{BiGUL}.TH}{}\<[E]%
\\
\>[B]{}\mathbf{import}\;\Conid{\Conid{Generics}.\Conid{BiGUL}.Lib}{}\<[E]%
\\
\>[B]{}\mathbf{import}\;\Conid{\Conid{Data}.List}{}\<[E]%
\\
\>[B]{}\mathbf{import}\;\Conid{\Conid{Data}.Maybe}{}\<[E]%
\\
\>[B]{}\mathbf{import}\;\Conid{\Conid{Control}.\Conid{Monad}.Except}{}\<[E]%
\\
\>[B]{}\mathbf{import}\;\Conid{\Conid{GHC}.Generics}{}\<[E]%
\ColumnHook
\end{hscode}\resethooks

}

Intuitively, we can think of a bidirectional BiGUL program:
\begin{hscode}\SaveRestoreHook
\column{B}{@{}>{\hspre}l<{\hspost}@{}}%
\column{3}{@{}>{\hspre}l<{\hspost}@{}}%
\column{E}{@{}>{\hspre}l<{\hspost}@{}}%
\>[3]{}\Varid{bx}\mathbin{::}\Conid{BiGUL}\;\Varid{s}\;\Varid{v}{}\<[E]%
\ColumnHook
\end{hscode}\resethooks
as describing how to manipulate a state consisting of
a source component of type~\ensuremath{\Varid{s}} and a view component of type \ensuremath{\Varid{v}};
the goal is to embed all information in the view to proper places
in the source. For each \ensuremath{\Varid{bx}\mathbin{::}\Conid{BiGUL}\;\Varid{s}\;\Varid{v}}, we can run it forwardly by calling \ensuremath{\Varid{get}}
and backwardly by calling \ensuremath{\Varid{put}}.
\begin{hscode}\SaveRestoreHook
\column{B}{@{}>{\hspre}l<{\hspost}@{}}%
\column{3}{@{}>{\hspre}l<{\hspost}@{}}%
\column{E}{@{}>{\hspre}l<{\hspost}@{}}%
\>[3]{}\Varid{get}\;\Varid{bx}\mathbin{::}\Varid{s}\to \Conid{Maybe}\;\Varid{v}{}\<[E]%
\\
\>[3]{}\Varid{put}\;\Varid{bx}\mathbin{::}\Varid{s}\to \Varid{v}\to \Conid{Maybe}\;\Varid{s}{}\<[E]%
\ColumnHook
\end{hscode}\resethooks
where \ensuremath{\Varid{get}\;\Varid{bx}} is a function mapping a source to a view if it succeeds,
while \ensuremath{\Varid{put}\;\Varid{bx}}
accepts an original source and uses a view to update it to get an updated source.

In BiGUL, it is suffice for users to write the \ensuremath{\Varid{put}} behavior (i.e.,
how to use a view to update the original source to a new source),
and the (unique) \ensuremath{\Varid{get}} behavior is obtained for free. 
%
The core of BiGUL consists of a small number of primitives and
combinators for constructing valid (well-behaved)
bidirectional transformations.

\subsection{Skip}
 
The first primitive for writing \ensuremath{\Varid{put}} is
\begin{hscode}\SaveRestoreHook
\column{B}{@{}>{\hspre}l<{\hspost}@{}}%
\column{3}{@{}>{\hspre}l<{\hspost}@{}}%
\column{12}{@{}>{\hspre}l<{\hspost}@{}}%
\column{E}{@{}>{\hspre}l<{\hspost}@{}}%
\>[3]{}\Conid{Skip}{}\<[12]%
\>[12]{}\mathbin{::}(\Varid{s}\to \Varid{v})\to \Conid{BiGUL}\;\Varid{s}\;\Varid{v}{}\<[E]%
\ColumnHook
\end{hscode}\resethooks
The put behavior of \ensuremath{\Conid{Skip}\;\Varid{f}} does not allow any change on the view
and thus skip any change on the source
(while in the get direction, the view is fully computed by
applying function \ensuremath{\Varid{f}} to the source). Consider a simple \ensuremath{\Varid{put}} defined by
\ensuremath{\Conid{Skip}\;\Varid{square}} where
\begin{hscode}\SaveRestoreHook
\column{B}{@{}>{\hspre}l<{\hspost}@{}}%
\column{E}{@{}>{\hspre}l<{\hspost}@{}}%
\>[B]{}\Varid{square}\;\Varid{x}\mathrel{=}\Varid{x}\mathbin{*}\Varid{x}{}\<[E]%
\ColumnHook
\end{hscode}\resethooks
and we can test it by
\begin{tabbing}\tt
~\char42{}Main\char62{}~put~\char40{}Skip~square\char41{}~10~100\\
\tt ~Just~10
\end{tabbing}
It first checks if the view \ensuremath{\mathrm{100}} is the square of the source \ensuremath{\mathrm{10}}, and returns
the original source if it is. But if the view is changed, say to \ensuremath{\mathrm{250}},
it should return \ensuremath{\Conid{Nothing}}:
\begin{tabbing}\tt
~\char42{}Main\char62{}~put~\char40{}Skip~square\char41{}~10~250\\
\tt ~Nothing
\end{tabbing}
If one wants to see why \ensuremath{\Varid{put}} returns \ensuremath{\Conid{Nothing}}, he may use
\ensuremath{\Varid{putTrace}} instead of \ensuremath{\Varid{put}} to get more information.
\begin{tabbing}\tt
~\char42{}Main\char62{}~putTrace~\char40{}Skip~square\char41{}~10~250\\
\tt ~view~not~determined~by~the~source
\end{tabbing}

Note that each well-behaved putback transformation in BiGUL
is equipped with a unique \ensuremath{\Varid{get}} for doing forward transformation.
We can test the \ensuremath{\Varid{get}} behavior as follows.
\begin{tabbing}\tt
~\char42{}Main\char62{}~get~\char40{}Skip~square\char41{}~5\\
\tt ~Just~25
\end{tabbing}
It reads that doing forward transformation for \ensuremath{\Conid{Skip}\;\Varid{square}} on the
source of \ensuremath{\mathrm{5}} gives the view of \ensuremath{\mathrm{25}}. Like that we have \ensuremath{\Varid{putTrace}} for \ensuremath{\Varid{put}}
to see more execution information about \ensuremath{\Varid{put}},
we have \ensuremath{\Varid{getTrace}} for \ensuremath{\Varid{get}} to see  more execution information about \ensuremath{\Varid{get}}.

As a simple exercise, can you see what the following \ensuremath{\Varid{skip1}} is? 
\begin{hscode}\SaveRestoreHook
\column{B}{@{}>{\hspre}l<{\hspost}@{}}%
\column{E}{@{}>{\hspre}l<{\hspost}@{}}%
\>[B]{}\Varid{skip1}\mathbin{::}\Conid{BiGUL}\;\Varid{s}\;(){}\<[E]%
\\
\>[B]{}\Varid{skip1}\mathrel{=}\Conid{Skip}\;(\Varid{const}\;()){}\<[E]%
\ColumnHook
\end{hscode}\resethooks

\subsection{Replace}

The second one is
\begin{hscode}\SaveRestoreHook
\column{B}{@{}>{\hspre}l<{\hspost}@{}}%
\column{3}{@{}>{\hspre}l<{\hspost}@{}}%
\column{12}{@{}>{\hspre}l<{\hspost}@{}}%
\column{E}{@{}>{\hspre}l<{\hspost}@{}}%
\>[3]{}\Conid{Replace}{}\<[12]%
\>[12]{}\mathbin{::}\Conid{BiGUL}\;\Varid{s}\;\Varid{s}{}\<[E]%
\ColumnHook
\end{hscode}\resethooks
to use the view to completely replace the source. For instance,
\begin{tabbing}\tt
~\char42{}Main\char62{}~put~Replace~1~100\\
\tt ~Just~100
\end{tabbing}
uses the view \ensuremath{\mathrm{100}} to replace the source \ensuremath{\mathrm{1}} and gets a new source \ensuremath{\mathrm{100}}.

\subsection{Product: Prod}

To use a view pair \ensuremath{(\Varid{v}_{1},\Varid{v}_{2})} to update a source pair
\ensuremath{(\Varid{s}_{1},\Varid{s}_{2})},
we can write \ensuremath{\Varid{bx1}\;\Conid{Prod}\;\Varid{bx2}}, a product of two putback transformations
\ensuremath{\Varid{bx1}} and \ensuremath{\Varid{bx2}},
which is to use \ensuremath{\Varid{v}_{1}} to update \ensuremath{\Varid{s}_{1}} with \ensuremath{\Varid{bx1}} and \ensuremath{\Varid{v}_{2}} to \ensuremath{\Varid{s}_{2}} with \ensuremath{\Varid{bx2}}.
\begin{hscode}\SaveRestoreHook
\column{B}{@{}>{\hspre}l<{\hspost}@{}}%
\column{3}{@{}>{\hspre}l<{\hspost}@{}}%
\column{E}{@{}>{\hspre}l<{\hspost}@{}}%
\>[3]{}\Conid{Prod}\mathbin{::}\Conid{BiGUL}\;\Varid{s}_{1}\;\Varid{v}_{1}\to \Conid{BiGUL}\;\Varid{s}_{2}\;\Varid{v}_{2}\to \Conid{BiGUL}\;(\Varid{s}_{1},\Varid{s}_{2})\;(\Varid{v}_{1},\Varid{v}_{2}){}\<[E]%
\ColumnHook
\end{hscode}\resethooks
For instance, we can combine \ensuremath{\Conid{Skip}} and \ensuremath{\Conid{Replace}} to putback a view pair
to a source pair.
\begin{tabbing}\tt
~\char42{}Main\char62{}~put~\char40{}skip1~\char96{}Prod\char96{}~Replace\char41{}~\char40{}5\char44{}1\char41{}~\char40{}\char40{}\char41{}\char44{}100\char41{}\\
\tt ~Just~\char40{}5\char44{}100\char41{}
\end{tabbing}
More complicated structures can be dealt with by using nested \ensuremath{\Conid{Prod}}:
\begin{tabbing}\tt
~\char42{}Main\char62{}~put~\char40{}\char40{}skip1~\char96{}Prod\char96{}~Replace\char41{}~\char96{}Prod\char96{}~Replace\char41{}\\
\tt ~~~~~~~~~~~~\char40{}\char40{}5\char44{}1\char41{}\char44{}2\char41{}~\char40{}\char40{}\char40{}\char41{}\char44{}100\char41{}\char44{}200\char41{}\\
\tt ~Just~\char40{}\char40{}5\char44{}100\char41{}\char44{}200\char41{}
\end{tabbing}

To help describe structural mapping more directly, we provide the following
syntactic sugar:\begin{hscode}\SaveRestoreHook
\column{B}{@{}>{\hspre}l<{\hspost}@{}}%
\column{3}{@{}>{\hspre}l<{\hspost}@{}}%
\column{E}{@{}>{\hspre}l<{\hspost}@{}}%
\>[3]{}\mathbin{\$}(\Varid{update}\;[\mskip1.5mu \Varid{p}\mid \Varid{source}\mathbin{-}\Varid{pattern}\mid \mskip1.5mu]\;[\mskip1.5mu \Varid{p}\mid \Varid{view}\mathbin{-}\Varid{pattern}\mid \mskip1.5mu]\;[\mskip1.5mu \Varid{d}\mid \Varid{updating}\mathbin{-}\Varid{strategy}\mid \mskip1.5mu]){}\<[E]%
\ColumnHook
\end{hscode}\resethooks
Source and view are decomposed by the patterns in the
\ensuremath{[\mskip1.5mu \Varid{p}} ... \ensuremath{\mskip1.5mu]} and corresponding elements are updated based on \ensuremath{\Varid{updating}\mathbin{-}\Varid{strategy}}.
As an example, we may describe \ensuremath{(\Varid{skip1}\mathbin{`\Conid{Prod}`}\Conid{Replace})\mathbin{`\Conid{Prod}`}\Conid{Replace}} by
\begin{hscode}\SaveRestoreHook
\column{B}{@{}>{\hspre}l<{\hspost}@{}}%
\column{24}{@{}>{\hspre}l<{\hspost}@{}}%
\column{E}{@{}>{\hspre}l<{\hspost}@{}}%
\>[B]{}\Varid{testUpdate}\mathbin{::}(\Conid{Show}\;\Varid{a},\Conid{Show}\;\Varid{b},\Conid{Show}\;\Varid{c})\Rightarrow \Conid{BiGUL}\;((\Varid{a},\Varid{b}),\Varid{c})\;(((),\Varid{b}),\Varid{c}){}\<[E]%
\\
\>[B]{}\Varid{testUpdate}\mathrel{=}\mathbin{\$}(\Varid{update}\;{}\<[24]%
\>[24]{}[\mskip1.5mu \Varid{p}\mid ((\Varid{x},\Varid{y}),\Varid{z})\mid \mskip1.5mu]\;[\mskip1.5mu \Varid{p}\mid ((\Varid{x},\Varid{y}),\Varid{z})\mid \mskip1.5mu]\;{}\<[E]%
\\
\>[24]{}[\mskip1.5mu \Varid{d}\mid \Varid{x}\mathrel{=}\Varid{skip1};\Varid{y}\mathrel{=}\Conid{Replace};\Varid{z}\mathrel{=}\Conid{Replace}\mid \mskip1.5mu]){}\<[E]%
\ColumnHook
\end{hscode}\resethooks
In this concrete example, the three elements of the tuple (both in the source and in the view) are bound to variables \ensuremath{\Varid{x},\Varid{y},\Varid{z}}, and they are sent to the three
combinators as arguments in the [d\ensuremath{\mathbin{...}}] part. Note that since \ensuremath{\Varid{skip1}} appears
quite often, we could just mark the source element as underline(\_) in the [p\ensuremath{\Varid{source}\mathbin{-}\Varid{pattern}}] and omit writing \ensuremath{\Varid{skip1}} in the [d\ensuremath{\mathbin{...}}] part.
\begin{hscode}\SaveRestoreHook
\column{B}{@{}>{\hspre}l<{\hspost}@{}}%
\column{25}{@{}>{\hspre}l<{\hspost}@{}}%
\column{E}{@{}>{\hspre}l<{\hspost}@{}}%
\>[B]{}\Varid{testUpdate'}\mathbin{::}(\Conid{Show}\;\Varid{a},\Conid{Show}\;\Varid{b},\Conid{Show}\;\Varid{c})\Rightarrow \Conid{BiGUL}\;((\Varid{a},\Varid{b}),\Varid{c})\;(((),\Varid{b}),\Varid{c}){}\<[E]%
\\
\>[B]{}\Varid{testUpdate'}\mathrel{=}\mathbin{\$}(\Varid{update}\;{}\<[25]%
\>[25]{}[\mskip1.5mu \Varid{p}\mid ((\anonymous ,\Varid{y}),\Varid{z})\mid \mskip1.5mu]\;[\mskip1.5mu \Varid{p}\mid (((),\Varid{y}),\Varid{z})\mid \mskip1.5mu]\;{}\<[E]%
\\
\>[25]{}[\mskip1.5mu \Varid{d}\mid \Varid{y}\mathrel{=}\Conid{Replace};\Varid{z}\mathrel{=}\Conid{Replace}\mid \mskip1.5mu]){}\<[E]%
\ColumnHook
\end{hscode}\resethooks

\subsection{Source/View rearrangement}

So far, the source and the view are of the same structure. What if we
wish to putback a view \ensuremath{(\Varid{v}_{1},\Varid{v}_{2})} to a source of a different structure,
say \ensuremath{((\Varid{s}_{0},\Varid{s}_{1}),\Varid{s}_{2}))}, to replace \ensuremath{\Varid{s}_{1}} by \ensuremath{\Varid{v}_{1}} and
\ensuremath{\Varid{s}_{2}} by \ensuremath{\Varid{v}_{2}}? BiGUL provides a way of 
rearranging either the source
or the view through a natural transformation \ensuremath{\Varid{tau}} to make both the view and
the source have the same structure.
\begin{hscode}\SaveRestoreHook
\column{B}{@{}>{\hspre}l<{\hspost}@{}}%
\column{3}{@{}>{\hspre}l<{\hspost}@{}}%
\column{E}{@{}>{\hspre}l<{\hspost}@{}}%
\>[3]{}\mathbin{\$}(\Varid{rearrS}\;[\mskip1.5mu \mid \Varid{tau}\mathbin{::}\Varid{s}_{1}\to \Varid{s}_{2}\mid \mskip1.5mu])\mathbin{::}\Conid{BiGUL}\;\Varid{s}_{2}\;\Varid{v}\to \Conid{BiGUL}\;\Varid{s}_{1}\;\Varid{v}{}\<[E]%
\\
\>[3]{}\mathbin{\$}(\Varid{rearrV}\;[\mskip1.5mu \mid \Varid{tau}\mathbin{::}\Varid{v}_{1}\to \Varid{v}_{2}\mid \mskip1.5mu])\mathbin{::}\Conid{BiGUL}\;\Varid{s}\;\Varid{v}_{2}\to \Conid{BiGUL}\;\Varid{s}\;\Varid{v}_{1}{}\<[E]%
\ColumnHook
\end{hscode}\resethooks
Returning to the above problem, we may define the following
putback transformation by rearranging the source
\begin{hscode}\SaveRestoreHook
\column{B}{@{}>{\hspre}l<{\hspost}@{}}%
\column{E}{@{}>{\hspre}l<{\hspost}@{}}%
\>[B]{}\Varid{putPairOverNPair}\mathbin{::}(\Conid{Show}\;\Varid{s}_{1},\Conid{Show}\;\Varid{s}_{2})\Rightarrow \Conid{BiGUL}\;((\Varid{s}_{0},\Varid{s}_{1}),\Varid{s}_{2})\;(\Varid{s}_{1},\Varid{s}_{2}){}\<[E]%
\\
\>[B]{}\Varid{putPairOverNPair}\mathrel{=}\mathbin{\$}(\Varid{rearrS}\;[\mskip1.5mu \mid \lambda ((\Varid{s}_{0},\Varid{s}_{1}),\Varid{s}_{2})\to (\Varid{s}_{1},\Varid{s}_{2})\mid \mskip1.5mu])\;\Conid{Replace}{}\<[E]%
\ColumnHook
\end{hscode}\resethooks
and have
\begin{tabbing}\tt
~\char42{}PBasic\char62{}~put~putPairOverNPair~\char40{}\char40{}5\char44{}1\char41{}\char44{}2\char41{}~\char40{}100\char44{}200\char41{}\\
\tt ~Just~\char40{}\char40{}5\char44{}100\char41{}\char44{}200\char41{}
\end{tabbing}
Or we may define it by rearranging the view:
\begin{hscode}\SaveRestoreHook
\column{B}{@{}>{\hspre}l<{\hspost}@{}}%
\column{22}{@{}>{\hspre}l<{\hspost}@{}}%
\column{24}{@{}>{\hspre}l<{\hspost}@{}}%
\column{E}{@{}>{\hspre}l<{\hspost}@{}}%
\>[B]{}\Varid{putPairOverNPair'}\mathbin{::}(\Conid{Show}\;\Varid{s}_{0},\Conid{Show}\;\Varid{s}_{1},\Conid{Show}\;\Varid{s}_{2})\Rightarrow \Conid{BiGUL}\;((\Varid{s}_{0},\Varid{s}_{1}),\Varid{s}_{2})\;(\Varid{s}_{1},\Varid{s}_{2}){}\<[E]%
\\
\>[B]{}\Varid{putPairOverNPair'}\mathrel{=}{}\<[22]%
\>[22]{}\mathbin{\$}(\Varid{rearrV}\;[\mskip1.5mu \mid \lambda (\Varid{v}_{1},\Varid{v}_{2})\to (((),\Varid{v}_{1}),\Varid{v}_{2})\mid \mskip1.5mu])\mathbin{\$}{}\<[E]%
\\
\>[22]{}\hsindent{2}{}\<[24]%
\>[24]{}(\Varid{skip1}\mathbin{`\Conid{Prod}`}\Conid{Replace})\mathbin{`\Conid{Prod}`}\Conid{Replace}{}\<[E]%
\ColumnHook
\end{hscode}\resethooks

The mechanism of source/view rearrangement enables us to
process algebraic data structures such as
lists and trees. The following example uses view \ensuremath{\Varid{v}} to
replace the first element of a nonempty source list.
\begin{hscode}\SaveRestoreHook
\column{B}{@{}>{\hspre}l<{\hspost}@{}}%
\column{E}{@{}>{\hspre}l<{\hspost}@{}}%
\>[B]{}\Varid{pHead}\mathbin{::}\Conid{Show}\;\Varid{s}\Rightarrow \Conid{BiGUL}\;[\mskip1.5mu \Varid{s}\mskip1.5mu]\;\Varid{s}{}\<[E]%
\\
\>[B]{}\Varid{pHead}\mathrel{=}\mathbin{\$}(\Varid{rearrS}\;[\mskip1.5mu \mid \lambda (\Varid{s}\mathbin{:\char95 })\to \Varid{s}\mid \mskip1.5mu])\;\Conid{Replace}{}\<[E]%
\ColumnHook
\end{hscode}\resethooks
\begin{tabbing}\tt
~\char42{}PBasic\char62{}~put~pHead~\char91{}1\char44{}2\char44{}3\char44{}4\char93{}~100\\
\tt ~Just~\char91{}100\char44{}2\char44{}3\char44{}4\char93{}
\end{tabbing}
What if we wish to define a general putback transformation 
that uses the view to replace the \ensuremath{\Varid{i}}th element of the source list?
We can define it recursively as follows.
\begin{hscode}\SaveRestoreHook
\column{B}{@{}>{\hspre}l<{\hspost}@{}}%
\column{11}{@{}>{\hspre}l<{\hspost}@{}}%
\column{17}{@{}>{\hspre}l<{\hspost}@{}}%
\column{20}{@{}>{\hspre}l<{\hspost}@{}}%
\column{23}{@{}>{\hspre}l<{\hspost}@{}}%
\column{E}{@{}>{\hspre}l<{\hspost}@{}}%
\>[B]{}\Varid{pNth}\mathbin{::}\Conid{Show}\;\Varid{s}\Rightarrow \mathit{Int}\to \Conid{BiGUL}\;[\mskip1.5mu \Varid{s}\mskip1.5mu]\;\Varid{s}{}\<[E]%
\\
\>[B]{}\Varid{pNth}\;\Varid{i}\mathrel{=}{}\<[11]%
\>[11]{}\mathbf{if}\;\Varid{i}\equiv \mathrm{0}\;\mathbf{then}\;\Varid{pHead}{}\<[E]%
\\
\>[11]{}\mathbf{else}{}\<[17]%
\>[17]{}\mathbin{\$}(\Varid{rearrS}\;[\mskip1.5mu \mid \lambda (\Varid{x}\mathbin{:}\Varid{xs})\to (\Varid{x},\Varid{xs})\mid \mskip1.5mu])\mathbin{\$}{}\<[E]%
\\
\>[17]{}\hsindent{3}{}\<[20]%
\>[20]{}\mathbin{\$}(\Varid{rearrV}\;[\mskip1.5mu \mid \lambda \Varid{v}\to ((),\Varid{v})\mid \mskip1.5mu])\mathbin{\$}{}\<[E]%
\\
\>[20]{}\hsindent{3}{}\<[23]%
\>[23]{}\Varid{skip1}\mathbin{`\Conid{Prod}`}\Varid{pNth}\;(\Varid{i}\mathbin{-}\mathrm{1}){}\<[E]%
\ColumnHook
\end{hscode}\resethooks
\begin{tabbing}\tt
~\char42{}PBasic\char62{}~put~\char40{}pNth~3\char41{}~\char91{}1\char46{}\char46{}10\char93{}~100\\
\tt ~Just~\char91{}1\char44{}2\char44{}3\char44{}100\char44{}5\char44{}6\char44{}7\char44{}8\char44{}9\char44{}10\char93{}
\end{tabbing}
As we know that any putback function in BiGUL is equipped with
a \ensuremath{\Varid{get}} function, for \ensuremath{\Varid{pNth}}, we can test its \ensuremath{\Varid{get}} behavior
as follows; its corresponding \ensuremath{\Varid{get}} function is actually
our familiar \ensuremath{\Varid{take}} function.
\begin{tabbing}\tt
~\char42{}PBasic\char62{}~get~\char40{}pNth~3\char41{}~\char91{}1\char46{}\char46{}10\char93{}\\
\tt ~Just~4
\end{tabbing}

\subsection{Case}

The \ensuremath{\Conid{Case}} combinator is for case analysis, which is very useful.
The general structure is as follows.\begin{hscode}\SaveRestoreHook
\column{B}{@{}>{\hspre}l<{\hspost}@{}}%
\column{3}{@{}>{\hspre}l<{\hspost}@{}}%
\column{6}{@{}>{\hspre}l<{\hspost}@{}}%
\column{9}{@{}>{\hspre}c<{\hspost}@{}}%
\column{9E}{@{}l@{}}%
\column{12}{@{}>{\hspre}l<{\hspost}@{}}%
\column{14}{@{}>{\hspre}l<{\hspost}@{}}%
\column{23}{@{}>{\hspre}l<{\hspost}@{}}%
\column{41}{@{}>{\hspre}l<{\hspost}@{}}%
\column{E}{@{}>{\hspre}l<{\hspost}@{}}%
\>[3]{}\Conid{Case}\;{}\<[9]%
\>[9]{}[\mskip1.5mu {}\<[9E]%
\>[12]{}\mathbin{\$}(\Varid{normal}\;{}\<[23]%
\>[23]{}[\mskip1.5mu \mid \Varid{enteringCond1}{}\<[41]%
\>[41]{}\mathbin{::}\Varid{s}\to \Varid{v}\to \Conid{Bool}\mid \mskip1.5mu]\;[\mskip1.5mu \mid \Varid{exitCond1}\mathbin{::}\Varid{s}\to \Conid{Bool}\mid \mskip1.5mu]){}\<[E]%
\\
\>[12]{}\hsindent{2}{}\<[14]%
\>[14]{}\Longrightarrow(\Varid{bx1}\mathbin{::}\Conid{BiGUL}\;\Varid{s}\;\Varid{v}){}\<[E]%
\\
\>[9]{},{}\<[9E]%
\>[12]{}\mathbin{\$}(\Varid{adaptive}\;[\mskip1.5mu \mid \Varid{enteringCond1'}\mathbin{::}\Varid{s}\to \Varid{v}\to \Conid{Bool}\mid \mskip1.5mu]){}\<[E]%
\\
\>[12]{}\hsindent{2}{}\<[14]%
\>[14]{}\Longrightarrow(\Varid{f1}\mathbin{::}\Varid{s}\to \Varid{v}\to \Varid{s}){}\<[E]%
\\
\>[9]{},{}\<[9E]%
\>[12]{}\mathbin{...}{}\<[E]%
\\
\>[9]{},{}\<[9E]%
\>[12]{}\mathbin{\$}(\Varid{normal}\;{}\<[23]%
\>[23]{}[\mskip1.5mu \mid \Varid{enteringCondn}{}\<[41]%
\>[41]{}\mathbin{::}\Varid{s}\to \Varid{v}\to \Conid{Bool}\mid \mskip1.5mu]\;[\mskip1.5mu \mid \Varid{exitCond1}\mathbin{::}\Varid{s}\to \Conid{Bool}\mid \mskip1.5mu]){}\<[E]%
\\
\>[12]{}\hsindent{2}{}\<[14]%
\>[14]{}\Longrightarrow(\Varid{bxn}\mathbin{::}\Conid{BiGUL}\;\Varid{s}\;\Varid{v}){}\<[E]%
\\
\>[9]{},{}\<[9E]%
\>[12]{}\mathbin{...}{}\<[E]%
\\
\>[9]{},{}\<[9E]%
\>[12]{}\mathbin{\$}(\Varid{adaptive}\;[\mskip1.5mu \mid \Varid{enteringCondm'}\mathbin{::}\Varid{s}\to \Varid{v}\to \Conid{Bool}\mid \mskip1.5mu]){}\<[E]%
\\
\>[12]{}\hsindent{2}{}\<[14]%
\>[14]{}\Longrightarrow(\Varid{fm}\mathbin{::}\Varid{s}\to \Varid{v}\to \Varid{s}){}\<[E]%
\\
\>[9]{}\mskip1.5mu]{}\<[9E]%
\\
\>[3]{}\hsindent{3}{}\<[6]%
\>[6]{}\mathbin{::}\Conid{BiGUL}\;\Varid{s}\;\Varid{v}{}\<[E]%
\ColumnHook
\end{hscode}\resethooks
It contains a sequence of cases. For each case, it is either normal or
adaptive. For the normal case, if the condition is satisfied, a
corresponding putback transformation is applied. For the adaptive
case, if the condition is satisfied, a function is used to update the
source with the view so that for the next step one of the normal cases
can be applied. Note that if adaptation does not lead the source and
the view to a normal case, an error will be reported at runtime.

Note that \ensuremath{\mathbin{\$}(\Varid{normal}\mathbin{...}\mathbin{...})} takes two predicates. The first one is
the entering-condition while the second one is the exit-condition. The
predicate for entering-condition is very general, and we can use any
function f of type \ensuremath{(\Varid{s}\to \Varid{v}\to \Conid{Bool})} to examine the source and view. If
the condition is matched, then the BiGUL program after the predicate
is executed. If the condition is not satisfied, the next branch is
tried. The predicate for exit-condition checks the source only. The
exit-condition in different branches should NOT be overlapped.

As a simple example, consider using the view to update all
the elements in the source list. To do so, we can use \ensuremath{\Conid{Case}}.
\begin{hscode}\SaveRestoreHook
\column{B}{@{}>{\hspre}l<{\hspost}@{}}%
\column{3}{@{}>{\hspre}l<{\hspost}@{}}%
\column{9}{@{}>{\hspre}c<{\hspost}@{}}%
\column{9E}{@{}l@{}}%
\column{12}{@{}>{\hspre}l<{\hspost}@{}}%
\column{14}{@{}>{\hspre}l<{\hspost}@{}}%
\column{21}{@{}>{\hspre}l<{\hspost}@{}}%
\column{24}{@{}>{\hspre}l<{\hspost}@{}}%
\column{E}{@{}>{\hspre}l<{\hspost}@{}}%
\>[B]{}\Varid{replaceAll}\mathbin{::}(\Conid{Eq}\;\Varid{s},\Conid{Show}\;\Varid{s})\Rightarrow \Conid{BiGUL}\;[\mskip1.5mu \Varid{s}\mskip1.5mu]\;\Varid{s}{}\<[E]%
\\
\>[B]{}\Varid{replaceAll}\mathrel{=}{}\<[E]%
\\
\>[B]{}\hsindent{3}{}\<[3]%
\>[3]{}\Conid{Case}\;{}\<[9]%
\>[9]{}[\mskip1.5mu {}\<[9E]%
\\
\>[9]{}\hsindent{3}{}\<[12]%
\>[12]{}\mathbin{\$}(\Varid{normal}\;[\mskip1.5mu \mid \lambda \Varid{s}\;\Varid{v}\to \Varid{length}\;\Varid{s}\equiv \mathrm{1}\mid \mskip1.5mu]\;[\mskip1.5mu \mid \lambda \Varid{s}\to \Varid{length}\;\Varid{s}\equiv \mathrm{1}\mid \mskip1.5mu]){}\<[E]%
\\
\>[12]{}\hsindent{2}{}\<[14]%
\>[14]{}\Longrightarrow\mathbin{\$}(\Varid{rearrS}\;[\mskip1.5mu \mid \lambda [\mskip1.5mu \Varid{x}\mskip1.5mu]\to \Varid{x}\mid \mskip1.5mu])\;\Conid{Replace},{}\<[E]%
\\
\>[9]{}\hsindent{3}{}\<[12]%
\>[12]{}\mathbin{\$}(\Varid{normal}\;[\mskip1.5mu \mid \lambda \Varid{s}\;\Varid{v}\to \Varid{length}\;\Varid{s}\mathbin{>}\mathrm{1}\mid \mskip1.5mu]\;[\mskip1.5mu \mid \lambda \Varid{s}\to \Varid{length}\;\Varid{s}\mathbin{>}\mathrm{1}\mid \mskip1.5mu]){}\<[E]%
\\
\>[12]{}\hsindent{2}{}\<[14]%
\>[14]{}\Longrightarrow\mathbin{\$}(\Varid{rearrS}\;[\mskip1.5mu \mid \lambda (\Varid{x}\mathbin{:}\Varid{xs})\to (\Varid{x},\Varid{xs})\mid \mskip1.5mu])\mathbin{\$}{}\<[E]%
\\
\>[14]{}\hsindent{7}{}\<[21]%
\>[21]{}\mathbin{\$}(\Varid{rearrV}\;[\mskip1.5mu \mid \lambda \Varid{v}\to (\Varid{v},\Varid{v})\mid \mskip1.5mu])\mathbin{\$}{}\<[E]%
\\
\>[21]{}\hsindent{3}{}\<[24]%
\>[24]{}\Conid{Replace}\mathbin{`\Conid{Prod}`}\Varid{replaceAll},{}\<[E]%
\\
\>[9]{}\hsindent{3}{}\<[12]%
\>[12]{}\mathbin{\$}(\Varid{adaptive}\;[\mskip1.5mu \mid \lambda \Varid{s}\;\Varid{v}\to \Varid{length}\;\Varid{s}\equiv \mathrm{0}\mid \mskip1.5mu]){}\<[E]%
\\
\>[12]{}\hsindent{2}{}\<[14]%
\>[14]{}\Longrightarrow\lambda \Varid{s}\;\Varid{v}\to [\mskip1.5mu \bot \mskip1.5mu]{}\<[E]%
\\
\>[9]{}\mskip1.5mu]{}\<[9E]%
\ColumnHook
\end{hscode}\resethooks
\begin{tabbing}\tt
~\char42{}PBasic\char62{}~put~replaceAll~\char91{}\char93{}~100\\
\tt ~Right~\char91{}100\char93{}\\
\tt ~\char42{}PBasic\char62{}~put~replaceAll~\char91{}1\char46{}\char46{}10\char93{}~100\\
\tt ~Right~\char91{}100\char44{}100\char44{}100\char44{}100\char44{}100\char44{}100\char44{}100\char44{}100\char44{}100\char44{}100\char93{}
\end{tabbing}

Note that instead of using a general function to describe
the entering condition or the exit condition, 
we can use patterns. The syntax for each case is:\begin{hscode}\SaveRestoreHook
\column{B}{@{}>{\hspre}l<{\hspost}@{}}%
\column{3}{@{}>{\hspre}l<{\hspost}@{}}%
\column{5}{@{}>{\hspre}l<{\hspost}@{}}%
\column{E}{@{}>{\hspre}l<{\hspost}@{}}%
\>[3]{}\mathbin{\$}(\Varid{normalSV}\;[\mskip1.5mu \Varid{p}\mid \Varid{source}\mathbin{-}\Varid{pattern}\mid \mskip1.5mu]\;[\mskip1.5mu \Varid{p}\mid \Varid{view}\mathbin{-}\Varid{pattern}\mid \mskip1.5mu]\;[\mskip1.5mu \mid \Varid{exitCond}\mid \mskip1.5mu]){}\<[E]%
\\
\>[3]{}\hsindent{2}{}\<[5]%
\>[5]{}\Longrightarrow(\Varid{bx}\mathbin{::}\Conid{BiGUL}\;\Varid{s}\;\Varid{v}){}\<[E]%
\\
\>[3]{}\mathbin{\$}(\Varid{adaptiveSV}\;[\mskip1.5mu \Varid{p}\mid \Varid{source}\mathbin{-}\Varid{pattern}\mid \mskip1.5mu]\;[\mskip1.5mu \Varid{p}\mid \Varid{view}\mathbin{-}\Varid{pattern}\mid \mskip1.5mu]){}\<[E]%
\\
\>[3]{}\hsindent{2}{}\<[5]%
\>[5]{}\Longrightarrow(\Varid{fm}\mathbin{::}\Varid{s}\to \Varid{v}\to \Varid{s}){}\<[E]%
\ColumnHook
\end{hscode}\resethooks
\ignore{
\begin{hscode}\SaveRestoreHook
\column{B}{@{}>{\hspre}l<{\hspost}@{}}%
\column{3}{@{}>{\hspre}l<{\hspost}@{}}%
\column{5}{@{}>{\hspre}l<{\hspost}@{}}%
\column{E}{@{}>{\hspre}l<{\hspost}@{}}%
\>[B]{}\Varid{repHead}\mathbin{::}\Conid{BiGUL}\;[\mskip1.5mu \mathit{Int}\mskip1.5mu]\;\mathit{Int}{}\<[E]%
\\
\>[B]{}\Varid{repHead}\mathrel{=}\Conid{Case}\;[\mskip1.5mu {}\<[E]%
\\
\>[B]{}\hsindent{3}{}\<[3]%
\>[3]{}\mathbin{\$}(\Varid{normal}\;[\mskip1.5mu \mid \lambda \Varid{s}\;\Varid{v}\to \Varid{length}\;\Varid{s}\mathbin{>}\mathrm{0}\mid \mskip1.5mu]\;[\mskip1.5mu \mid \lambda \Varid{s}\to \Varid{length}\;\Varid{s}\mathbin{>}\mathrm{0}\mid \mskip1.5mu]){}\<[E]%
\\
\>[3]{}\hsindent{2}{}\<[5]%
\>[5]{}\Longrightarrow\mathbin{\$}(\Varid{rearrS}\;[\mskip1.5mu \mid \lambda (\Varid{x}\mathbin{:}\Varid{xs})\to \Varid{x}\mid \mskip1.5mu])\;\Conid{Replace},{}\<[E]%
\\
\>[B]{}\hsindent{3}{}\<[3]%
\>[3]{}\mathbin{\$}(\Varid{adaptive}\;[\mskip1.5mu \mid \lambda \Varid{s}\;\Varid{v}\to \Varid{length}\;\Varid{s}\equiv \mathrm{0}\mid \mskip1.5mu]){}\<[E]%
\\
\>[3]{}\hsindent{2}{}\<[5]%
\>[5]{}\Longrightarrow\lambda \Varid{s}\;\Varid{v}\to [\mskip1.5mu \mathrm{0}\mskip1.5mu]{}\<[E]%
\\
\>[B]{}\mskip1.5mu]{}\<[E]%
\ColumnHook
\end{hscode}\resethooks
}

\subsection{Dependency: Dep}

Sometimes, a view may contain derived value that is computed from
other part of the view and not allowed
to be changed. For instance, for the view \ensuremath{(\Varid{x},\Varid{x}\mathbin{+}\mathrm{1})}, the second
component can be computed from the first by increasing it by \ensuremath{\mathrm{1}}.
To capture this, BiGUL provides
\begin{hscode}\SaveRestoreHook
\column{B}{@{}>{\hspre}l<{\hspost}@{}}%
\column{3}{@{}>{\hspre}l<{\hspost}@{}}%
\column{E}{@{}>{\hspre}l<{\hspost}@{}}%
\>[3]{}\Conid{Dep}\mathbin{::}\Conid{Eq}\;\Varid{v'}\Rightarrow (\Varid{v}\to \Varid{v'})\to \Conid{BiGUL}\;\Varid{a}\;\Varid{v}\to \Conid{BiGUL}\;\Varid{a}\;(\Varid{v},\Varid{v'}){}\<[E]%
\ColumnHook
\end{hscode}\resethooks
to describe this intention. We may, for example, define
\begin{hscode}\SaveRestoreHook
\column{B}{@{}>{\hspre}l<{\hspost}@{}}%
\column{E}{@{}>{\hspre}l<{\hspost}@{}}%
\>[B]{}\Varid{replaceAll2}\mathbin{::}\Conid{BiGUL}\;[\mskip1.5mu \mathit{Int}\mskip1.5mu]\;(\mathit{Int},\mathit{Int}){}\<[E]%
\\
\>[B]{}\Varid{replaceAll2}\mathrel{=}\Conid{Dep}\;(\mathbin{+}\mathrm{1})\;\Varid{replaceAll}{}\<[E]%
\ColumnHook
\end{hscode}\resethooks
to replace all elements of the source by the
first component of the view, while the second component
of the view is a derived one that should be one bigger than the first one.
\begin{tabbing}\tt
~\char42{}PBasic\char62{}~put~replaceAll2~\char91{}1\char46{}\char46{}10\char93{}~\char40{}100\char44{}101\char41{}\\
\tt ~Just~\char91{}100\char44{}100\char44{}100\char44{}100\char44{}100\char44{}100\char44{}100\char44{}100\char44{}100\char44{}100\char93{}\\
\tt ~\char42{}PBasic\char62{}~put~replaceAll2~\char91{}1\char46{}\char46{}10\char93{}~\char40{}100\char44{}200\char41{}\\
\tt ~Nothing\\
\tt ~\char42{}PBasic\char62{}~putTrace~replaceAll2~\char91{}1\char46{}\char46{}10\char93{}~\char40{}100\char44{}200\char41{}\\
\tt ~second~view~component~not~determined~by~the~first
\end{tabbing}
As seen in the last running of \ensuremath{\Varid{put}}, it reports an error because in the view \ensuremath{(\mathrm{100},\mathrm{200})},
\ensuremath{\mathrm{200}} is not one bigger than \ensuremath{\mathrm{100}}.

\subsection{Composition}

Putback transformations can be composed.
\begin{hscode}\SaveRestoreHook
\column{B}{@{}>{\hspre}l<{\hspost}@{}}%
\column{3}{@{}>{\hspre}l<{\hspost}@{}}%
\column{E}{@{}>{\hspre}l<{\hspost}@{}}%
\>[3]{}\Conid{Compose}\mathbin{::}\Conid{BiGUL}\;\Varid{a}\;\Varid{u}\to \Conid{BiGUL}\;\Varid{u}\;\Varid{b}\to \Conid{BiGUL}\;\Varid{a}\;\Varid{b}{}\<[E]%
\ColumnHook
\end{hscode}\resethooks
As a simple example, consider the following definition.
\begin{hscode}\SaveRestoreHook
\column{B}{@{}>{\hspre}l<{\hspost}@{}}%
\column{E}{@{}>{\hspre}l<{\hspost}@{}}%
\>[B]{}\Varid{pHead2}\mathbin{::}\Conid{Show}\;\Varid{a}\Rightarrow \Conid{BiGUL}\;[\mskip1.5mu [\mskip1.5mu \Varid{a}\mskip1.5mu]\mskip1.5mu]\;\Varid{a}{}\<[E]%
\\
\>[B]{}\Varid{pHead2}\mathrel{=}\Varid{pHead}\mathbin{`\Conid{Compose}`}\Varid{pHead}{}\<[E]%
\ColumnHook
\end{hscode}\resethooks
\begin{tabbing}\tt
~\char42{}PBasic\char62{}~put~pHead2~\char91{}\char91{}1\char44{}2\char93{}\char44{}\char91{}3\char44{}4\char44{}5\char93{}\char44{}\char91{}\char93{}\char93{}~100\\
\tt ~Just~\char91{}\char91{}100\char44{}2\char93{}\char44{}\char91{}3\char44{}4\char44{}5\char93{}\char44{}\char91{}\char93{}\char93{}
\end{tabbing}

\subsection{Utilities}

Generics.BiGUL.Lib has predefined some useful functions for building putback transformations.
One interesting one is \ensuremath{\Varid{emb}} that can safely embed a pair of well-behaved
\ensuremath{\Varid{get}} and \ensuremath{\Varid{put}} a putback transformation into our context.
\ensuremath{\Varid{emb}} itself is defined as follows.
\begin{hscode}\SaveRestoreHook
\column{B}{@{}>{\hspre}l<{\hspost}@{}}%
\column{3}{@{}>{\hspre}l<{\hspost}@{}}%
\column{5}{@{}>{\hspre}l<{\hspost}@{}}%
\column{7}{@{}>{\hspre}l<{\hspost}@{}}%
\column{E}{@{}>{\hspre}l<{\hspost}@{}}%
\>[3]{}\Varid{emb}\mathbin{::}\Conid{Eq}\;\Varid{v}\Rightarrow (\Varid{s}\to \Varid{v})\to (\Varid{s}\to \Varid{v}\to \Varid{s})\to \Conid{BiGUL}\;\Varid{s}\;\Varid{v}{}\<[E]%
\\
\>[3]{}\Varid{emb}\;\Varid{g}\;\Varid{p}\mathrel{=}\Conid{Case}{}\<[E]%
\\
\>[3]{}\hsindent{2}{}\<[5]%
\>[5]{}[\mskip1.5mu \mathbin{\$}(\Varid{normal}\;[\mskip1.5mu \mid \lambda \Varid{s}\;\Varid{v}\to \Varid{g}\;\Varid{s}\equiv \Varid{v}\mid \mskip1.5mu]\;[\mskip1.5mu \Varid{p}\mid \anonymous \mid \mskip1.5mu]){}\<[E]%
\\
\>[5]{}\hsindent{2}{}\<[7]%
\>[7]{}\Longrightarrow\Conid{Skip}\;\Varid{g}{}\<[E]%
\\
\>[3]{}\hsindent{2}{}\<[5]%
\>[5]{},\mathbin{\$}(\Varid{adaptive}\;[\mskip1.5mu \mid \lambda \Varid{s}\;\Varid{v}\to \mbox{\commentbegin  g s /= v  \commentend}\Conid{True}\mid \mskip1.5mu]){}\<[E]%
\\
\>[5]{}\hsindent{2}{}\<[7]%
\>[7]{}\Longrightarrow\Varid{p}{}\<[E]%
\\
\>[3]{}\hsindent{2}{}\<[5]%
\>[5]{}\mskip1.5mu]{}\<[E]%
\ColumnHook
\end{hscode}\resethooks
As an application of \ensuremath{\Varid{emb}}, we may define the following putback function
for using the sum to update a pair.
\begin{hscode}\SaveRestoreHook
\column{B}{@{}>{\hspre}l<{\hspost}@{}}%
\column{3}{@{}>{\hspre}l<{\hspost}@{}}%
\column{10}{@{}>{\hspre}l<{\hspost}@{}}%
\column{E}{@{}>{\hspre}l<{\hspost}@{}}%
\>[B]{}\Varid{distSum}\mathbin{::}\Conid{BiGUL}\;(\mathit{Int},\mathit{Int})\;\mathit{Int}{}\<[E]%
\\
\>[B]{}\Varid{distSum}\mathrel{=}\Varid{emb}\;\Varid{g}\;\Varid{p}{}\<[E]%
\\
\>[B]{}\hsindent{3}{}\<[3]%
\>[3]{}\mathbf{where}\;{}\<[10]%
\>[10]{}\Varid{g}\;(\Varid{x},\Varid{y})\mathrel{=}\Varid{x}\mathbin{+}\Varid{y}{}\<[E]%
\\
\>[10]{}\Varid{p}\;(\Varid{x},\Varid{y})\;\Varid{v}\mathrel{=}(\Varid{v}\mathbin{-}\Varid{y},\Varid{y}){}\<[E]%
\ColumnHook
\end{hscode}\resethooks
\begin{tabbing}\tt
~\char42{}PBasic\char62{}~put~distSum~\char40{}1\char44{}2\char41{}~100\\
\tt ~Right~\char40{}98\char44{}2\char41{}\\
\tt ~\char42{}PBasic\char62{}~get~distSum~\char40{}1\char44{}2\char41{}\\
\tt ~Right~3
\end{tabbing}

\ignore{
\begin{hscode}\SaveRestoreHook
\column{B}{@{}>{\hspre}l<{\hspost}@{}}%
\column{3}{@{}>{\hspre}l<{\hspost}@{}}%
\column{5}{@{}>{\hspre}l<{\hspost}@{}}%
\column{6}{@{}>{\hspre}l<{\hspost}@{}}%
\column{8}{@{}>{\hspre}l<{\hspost}@{}}%
\column{9}{@{}>{\hspre}l<{\hspost}@{}}%
\column{E}{@{}>{\hspre}l<{\hspost}@{}}%
\>[B]{}\Varid{pHead'}\mathbin{::}\Conid{Show}\;\Varid{s}{}\<[E]%
\\
\>[B]{}\hsindent{8}{}\<[8]%
\>[8]{}\Rightarrow \Conid{BiGUL}\;[\mskip1.5mu \Varid{s}\mskip1.5mu]\;\Varid{s}{}\<[E]%
\\
\>[B]{}\Varid{pHead'}\mathrel{=}\Conid{Case}\;[\mskip1.5mu {}\<[E]%
\\
\>[B]{}\hsindent{6}{}\<[6]%
\>[6]{}\mathbin{\$}(\Varid{normal}\;[\mskip1.5mu \mid \lambda \Varid{s}\;\Varid{v}\to not\;(\Varid{null}\;\Varid{s})\mid \mskip1.5mu]\;[\mskip1.5mu \mid not\mathbin{\circ}\Varid{null}\mid \mskip1.5mu]){}\<[E]%
\\
\>[6]{}\hsindent{2}{}\<[8]%
\>[8]{}\Longrightarrow\Varid{pHead},{}\<[E]%
\\
\>[B]{}\hsindent{6}{}\<[6]%
\>[6]{}\mathbin{\$}(\Varid{adaptive}\;[\mskip1.5mu \mid \lambda \Varid{s}\;\Varid{v}\to \Varid{null}\;\Varid{s}\mid \mskip1.5mu]){}\<[E]%
\\
\>[6]{}\hsindent{2}{}\<[8]%
\>[8]{}\Longrightarrow\lambda \Varid{s}\;\Varid{v}\to [\mskip1.5mu \Varid{v}\mskip1.5mu]{}\<[E]%
\\
\>[B]{}\hsindent{6}{}\<[6]%
\>[6]{}\mskip1.5mu]{}\<[E]%
\\[\blanklineskip]%
\>[B]{}\Varid{pEither}\mathbin{::}(\Conid{Show}\;\Varid{a},\Conid{Show}\;\Varid{b},\Conid{Eq}\;\Varid{a}){}\<[E]%
\\
\>[B]{}\hsindent{9}{}\<[9]%
\>[9]{}\Rightarrow \Varid{a}\to \Conid{BiGUL}\;(\Conid{Either}\;\Varid{a}\;\Varid{b})\;\Varid{a}{}\<[E]%
\\
\>[B]{}\Varid{pEither}\;\Varid{x0}\mathrel{=}\Conid{Case}\;[\mskip1.5mu {}\<[E]%
\\
\>[B]{}\hsindent{3}{}\<[3]%
\>[3]{}\mathbin{\$}(\Varid{normalSV}\;[\mskip1.5mu \Varid{p}\mid \Conid{Left}\;\anonymous \mid \mskip1.5mu]\;[\mskip1.5mu \Varid{p}\mid \anonymous \mid \mskip1.5mu]\;[\mskip1.5mu \Varid{p}\mid \Conid{Left}\;\anonymous \mid \mskip1.5mu]){}\<[E]%
\\
\>[3]{}\hsindent{2}{}\<[5]%
\>[5]{}\Longrightarrow\mathbin{\$}(\Varid{update}\;[\mskip1.5mu \Varid{p}\mid \Conid{Left}\;\Varid{x}\mid \mskip1.5mu]\;[\mskip1.5mu \Varid{p}\mid \Varid{x}\mid \mskip1.5mu]\;[\mskip1.5mu \Varid{d}\mid \Varid{x}\mathrel{=}\Conid{Replace}\mid \mskip1.5mu]),{}\<[E]%
\\
\>[B]{}\hsindent{3}{}\<[3]%
\>[3]{}\mathbin{\$}(\Varid{normalSV}\;[\mskip1.5mu \Varid{p}\mid \Conid{Right}\;\Varid{y}\mid \mskip1.5mu]\;[\mskip1.5mu \mid \lambda \Varid{x}\to \Varid{x}\equiv \Varid{x0}\mid \mskip1.5mu]\;[\mskip1.5mu \Varid{p}\mid \Conid{Right}\;\anonymous \mid \mskip1.5mu]){}\<[E]%
\\
\>[3]{}\hsindent{2}{}\<[5]%
\>[5]{}\Longrightarrow\Conid{Skip}\;(\Varid{const}\;\Varid{x0}),{}\<[E]%
\\
\>[B]{}\hsindent{3}{}\<[3]%
\>[3]{}\mathbin{\$}(\Varid{adaptive}\;[\mskip1.5mu \mid \lambda \anonymous \;\anonymous \to \Conid{True}\mid \mskip1.5mu]){}\<[E]%
\\
\>[3]{}\hsindent{2}{}\<[5]%
\>[5]{}\Longrightarrow\lambda \Varid{s}\;\Varid{v}\to \Conid{Left}\;\Varid{v}{}\<[E]%
\\
\>[B]{}\hsindent{3}{}\<[3]%
\>[3]{}\mskip1.5mu]{}\<[E]%
\ColumnHook
\end{hscode}\resethooks
}
